\documentclass[12pt,a4paper]{article}

\usepackage[T2A]{fontenc}
\usepackage[utf8]{inputenc}
\usepackage[russian]{babel}
\usepackage{amsmath,amssymb,amsthm}
\usepackage{mathtools}
\usepackage{geometry}
\usepackage{array}
\usepackage{booktabs}
\usepackage{listings}
\usepackage{xcolor}
\usepackage{tikz}
\usepackage{hyperref}

\geometry{margin=2cm}

\lstdefinestyle{cuda}{
    language=C++,
    basicstyle=\ttfamily\small,
    keywordstyle=\color{blue},
    commentstyle=\color{gray},
    stringstyle=\color{teal},
    numbers=left,
    numberstyle=\tiny,
    stepnumber=1,
    frame=single,
    breaklines=true,
    tabsize=4
}

\newtheorem{definition}{Определение}
\newtheorem{theorem}{Теорема}
\newtheorem{proposition}{Утверждение}
\newtheorem{remark}{Замечание}

\title{Билет №28 \\ \small Графические процессоры и их применение для решения вычислительных задач. (ИТМО)}
\author{Сериков Владислав Александрович}
\date{6 мая 2026}

\begin{document}

\maketitle
\tableofcontents
\newpage

\section{Общее представление о графических процессорах}

\begin{definition}
\textbf{Графический процессор} --- это специализированный многопроцессорный вычислительный ускоритель, изначально предназначенный для обработки графики, но в современной вычислительной технике широко применяемый для массово-параллельных вычислений общего назначения.
\end{definition}

Традиционно центральный процессор, или CPU, оптимизируется для быстрого выполнения относительно небольшого числа сложных последовательных потоков управления. Графический процессор, или GPU, напротив, оптимизируется для выполнения большого числа простых однотипных операций над большими массивами данных.

Основная идея применения GPU в вычислительных задачах состоит в том, чтобы перенести на GPU те части программы, которые допускают массовый параллелизм:
\[
    y_i = f(x_i), \qquad i = 0,1,\dots,n-1.
\]
Если значения $y_i$ можно вычислять независимо или почти независимо, то такая задача хорошо подходит для GPU.

\subsection{CPU и GPU: различие архитектурных акцентов}

\begin{center}
\begin{tabular}{p{0.25\textwidth}p{0.32\textwidth}p{0.32\textwidth}}
\toprule
\textbf{Характеристика} & \textbf{CPU} & \textbf{GPU} \\
\midrule
Основная цель & Минимизация задержки одного потока & Максимизация пропускной способности \\
Число ядер & Обычно десятки & Сотни, тысячи арифметических устройств \\
Кэш-память & Большая и сложная иерархия кэшей & Меньше кэша на поток, но высокая пропускная способность памяти \\
Управление & Сложное предсказание ветвлений, внеочередное выполнение & Простые управляющие устройства, много потоков \\
Типичные задачи & ОС, компиляторы, базы данных, сложная логика & Линейная алгебра, обработка изображений, нейросети, моделирование \\
\bottomrule
\end{tabular}
\end{center}

GPU эффективен не потому, что каждый отдельный поток выполняется быстрее, а потому, что одновременно выполняется очень много потоков. Поэтому GPU ориентирован на \textbf{throughput computing} --- вычисления с высокой суммарной пропускной способностью.

\section{GPGPU: вычисления общего назначения на GPU}

\begin{definition}
\textbf{GPGPU} от англ. \emph{General-Purpose computing on Graphics Processing Units} --- использование графических процессоров для неграфических вычислительных задач общего назначения.
\end{definition}

Примеры задач GPGPU:

\begin{itemize}
    \item умножение матриц;
    \item решение систем линейных уравнений;
    \item численное решение дифференциальных уравнений;
    \item молекулярная динамика;
    \item обработка изображений и видео;
    \item обучение и инференс нейронных сетей;
    \item трассировка лучей;
    \item Монте-Карло моделирование;
    \item криптографические и переборные задачи;
    \item вычислительная гидродинамика.
\end{itemize}

\section{Модель исполнения GPU}

\subsection{Хост и устройство}

При программировании GPU обычно различают:

\begin{itemize}
    \item \textbf{host} --- центральный процессор и основная программа;
    \item \textbf{device} --- графический процессор, на котором запускаются параллельные ядра.
\end{itemize}

Типичная схема работы:

\begin{enumerate}
    \item CPU выделяет память на GPU.
    \item CPU копирует входные данные из оперативной памяти в память GPU.
    \item CPU запускает на GPU специальную функцию --- \textbf{ядро}.
    \item GPU выполняет ядро множеством потоков.
    \item CPU копирует результат из памяти GPU обратно.
\end{enumerate}

\subsection{Ядро}

\begin{definition}
\textbf{Ядро} от англ. \emph{kernel} --- функция, исполняемая на GPU большим числом параллельных потоков.
\end{definition}

Каждый поток обычно выполняет один и тот же код, но работает со своим индексом данных. Например, при сложении векторов поток с номером $i$ вычисляет
\[
    c_i = a_i + b_i.
\]

\section{Модель SIMT}

\begin{definition}
\textbf{SIMT} от англ. \emph{Single Instruction, Multiple Threads} --- модель исполнения, при которой множество потоков выполняет одну и ту же инструкцию над разными данными, но каждый поток имеет собственные регистры, индекс и состояние.
\end{definition}

SIMT близка к SIMD, но более гибкая. В SIMD одна инструкция применяется к вектору данных как единая операция. В SIMT программист мыслит отдельными потоками, а аппаратная часть группирует их для совместного исполнения.

\subsection{Варпы}

В архитектурах NVIDIA базовой группой потоков является \textbf{warp}. Обычно warp содержит $32$ потока.

\begin{definition}
\textbf{Warp} --- группа потоков GPU, которые планировщик исполняет совместно по модели SIMT.
\end{definition}

Если все потоки warp выполняют одну и ту же инструкцию, GPU работает эффективно. Если потоки расходятся по разным ветвям условного оператора, возникает \textbf{дивергенция}.

\subsection{Дивергенция потоков}

Рассмотрим код:

\begin{lstlisting}[style=cuda]
if (threadIdx.x % 2 == 0) {
    x = x + 1;
} else {
    x = x - 1;
}
\end{lstlisting}

Половина потоков warp идет в ветку \texttt{if}, половина --- в ветку \texttt{else}. GPU не может выполнить обе ветки одновременно для одного warp. Поэтому он последовательно выполняет одну ветку с активной маской одних потоков, затем другую ветку с активной маской остальных потоков.

\[
\begin{array}{c|cccccccc}
\text{Поток} & 0 & 1 & 2 & 3 & 4 & 5 & 6 & 7 \\
\hline
\text{Ветка} & A & B & A & B & A & B & A & B
\end{array}
\]

Эффективность падает, потому что часть потоков временно простаивает.

\begin{remark}
Дивергенция особенно вредна, если она возникает внутри одного warp. Если разные блоки потоков выполняют разные ветви, это обычно менее критично.
\end{remark}

\section{Иерархия потоков CUDA}

В CUDA потоки организуются в иерархию:

\[
    \text{grid} \supset \text{blocks} \supset \text{threads}.
\]

\begin{itemize}
    \item \textbf{thread} --- отдельный поток исполнения;
    \item \textbf{block} --- группа потоков, которые могут взаимодействовать через быструю разделяемую память;
    \item \textbf{grid} --- совокупность блоков, запускаемых одним вызовом ядра.
\end{itemize}

\subsection{Одномерная индексация}

Если используется одномерная сетка одномерных блоков, глобальный индекс потока равен:
\[
    i = \texttt{blockIdx.x} \cdot \texttt{blockDim.x} + \texttt{threadIdx.x}.
\]

\subsection{Двумерная индексация}

Для двумерной задачи, например изображения или матрицы:
\[
    x = \texttt{blockIdx.x} \cdot \texttt{blockDim.x} + \texttt{threadIdx.x},
\]
\[
    y = \texttt{blockIdx.y} \cdot \texttt{blockDim.y} + \texttt{threadIdx.y}.
\]

Если матрица хранится в строковом порядке, то линейный индекс элемента $(y,x)$ равен:
\[
    i = y \cdot W + x,
\]
где $W$ --- число столбцов.

\section{Иерархия памяти GPU}

Производительность GPU-программы часто определяется не количеством арифметических операций, а эффективностью доступа к памяти.

\subsection{Основные виды памяти}

\begin{center}
\begin{tabular}{p{0.23\textwidth}p{0.32\textwidth}p{0.32\textwidth}}
\toprule
\textbf{Вид памяти} & \textbf{Особенности} & \textbf{Типичное применение} \\
\midrule
Регистры & Самая быстрая память, индивидуальна для потока & Локальные переменные \\
Разделяемая память & Быстрая память внутри блока & Обмен данными между потоками блока, тайлинг \\
Глобальная память & Большая, но сравнительно медленная & Основные входные и выходные массивы \\
Константная память & Кэшируемая память для неизменяемых данных & Константы, параметры \\
Текстурная память & Оптимизирована для пространственной локальности & Обработка изображений \\
Локальная память & Формально локальна для потока, но физически может размещаться в глобальной памяти & Большие локальные массивы, spill регистров \\
\bottomrule
\end{tabular}
\end{center}

\subsection{Схема иерархии памяти}

\[
\begin{array}{c}
\boxed{\text{Регистры потока}} \\
\downarrow \\
\boxed{\text{Разделяемая память блока}} \\
\downarrow \\
\boxed{\text{L1/L2-кэши}} \\
\downarrow \\
\boxed{\text{Глобальная память GPU}} \\
\downarrow \\
\boxed{\text{Память CPU}}
\end{array}
\]

Чем ближе память к вычислительным устройствам, тем она быстрее, но тем меньше ее объем.

\section{Коалесцированный доступ к памяти}

\begin{definition}
\textbf{Коалесцированный доступ} --- такой шаблон обращения потоков warp к глобальной памяти, при котором обращения соседних потоков идут к соседним адресам и объединяются аппаратурой в малое число транзакций памяти.
\end{definition}

Пусть поток $t$ обращается к элементу массива $a[t]$. Тогда потоки warp обращаются к последовательным адресам:
\[
    a[0], a[1], a[2], \dots, a[31].
\]
Такой доступ обычно эффективен.

Если же поток $t$ обращается к $a[32t]$, адреса сильно разнесены:
\[
    a[0], a[32], a[64], \dots,
\]
и GPU вынужден выполнить больше транзакций памяти.

\subsection{Минимальный пример}

Хороший доступ:

\begin{lstlisting}[style=cuda]
int i = blockIdx.x * blockDim.x + threadIdx.x;
c[i] = a[i] + b[i];
\end{lstlisting}

Плохой доступ:

\begin{lstlisting}[style=cuda]
int i = blockIdx.x * blockDim.x + threadIdx.x;
c[i] = a[32 * i] + b[32 * i];
\end{lstlisting}

\section{Разделяемая память и тайлинг}

\begin{definition}
\textbf{Тайлинг} --- разбиение задачи на небольшие блоки данных, которые загружаются в быструю локальную память и многократно используются.
\end{definition}

Тайлинг особенно важен при умножении матриц. Пусть нужно вычислить:
\[
    C = AB,
\]
где
\[
    C_{ij} = \sum_{k=0}^{n-1} A_{ik}B_{kj}.
\]

В наивном варианте каждый поток многократно читает элементы из глобальной памяти. При тайлинге блок потоков загружает подматрицы $A$ и $B$ в разделяемую память, затем выполняет несколько операций умножения-сложения, используя быстрый доступ.

\subsection{Иллюстрация тайлинга}

Пусть матрицы разбиваются на плитки размера $T \times T$:

\[
A =
\begin{bmatrix}
A_{00} & A_{01} & \cdots \\
A_{10} & A_{11} & \cdots \\
\vdots & \vdots & \ddots
\end{bmatrix},
\qquad
B =
\begin{bmatrix}
B_{00} & B_{01} & \cdots \\
B_{10} & B_{11} & \cdots \\
\vdots & \vdots & \ddots
\end{bmatrix}.
\]

Для вычисления плитки $C_{ij}$ нужно:
\[
    C_{ij} = \sum_k A_{ik}B_{kj}.
\]

\section{Основные классы задач, подходящих для GPU}

Задача хорошо подходит для GPU, если выполняются следующие условия:

\begin{enumerate}
    \item Большой объем однотипных вычислений.
    \item Высокая степень параллелизма.
    \item Небольшое число зависимостей между элементами данных.
    \item Регулярный доступ к памяти.
    \item Высокая арифметическая интенсивность.
\end{enumerate}

\begin{definition}
\textbf{Арифметическая интенсивность} --- отношение числа арифметических операций к объему переданных данных:
\[
    I = \frac{\text{число операций}}{\text{число байт, переданных из памяти}}.
\]
\end{definition}

Чем выше арифметическая интенсивность, тем больше вероятность, что задача будет эффективно выполняться на GPU.

\section{Закон Амдала}

При оценке ускорения от использования GPU важно учитывать, что не вся программа может быть распараллелена.

\begin{theorem}[Закон Амдала]
Пусть доля времени выполнения программы, допускающая ускорение, равна $p$, где $0 \le p \le 1$. Пусть эта часть ускоряется в $s$ раз, где $s \ge 1$. Тогда общее ускорение программы равно
\[
    S = \frac{1}{(1-p) + \frac{p}{s}}.
\]
При $s \to \infty$ максимальное ускорение ограничено величиной
\[
    S_{\max} = \frac{1}{1-p}.
\]
\end{theorem}

\begin{proof}
Пусть исходное время выполнения программы равно $T$. Тогда последовательная часть занимает
\[
    T_{\text{seq}} = (1-p)T,
\]
а ускоряемая часть занимает
\[
    T_{\text{par}} = pT.
\]

После ускорения в $s$ раз время ускоряемой части станет:
\[
    T'_{\text{par}} = \frac{pT}{s}.
\]

Новое полное время:
\[
    T' = (1-p)T + \frac{pT}{s}
       = T\left((1-p)+\frac{p}{s}\right).
\]

Ускорение определяется как отношение старого времени к новому:
\[
    S = \frac{T}{T'}
      = \frac{T}{T\left((1-p)+\frac{p}{s}\right)}
      = \frac{1}{(1-p)+\frac{p}{s}}.
\]

Если $s \to \infty$, то
\[
    \frac{p}{s} \to 0,
\]
поэтому
\[
    S_{\max} = \frac{1}{1-p}.
\]
Теорема доказана.
\end{proof}

\subsection{Пример применения закона Амдала}

Пусть $90\%$ программы можно ускорить на GPU в $100$ раз. Тогда:
\[
    p = 0.9,\qquad s = 100.
\]
Получаем:
\[
    S = \frac{1}{0.1 + \frac{0.9}{100}}
      = \frac{1}{0.1 + 0.009}
      = \frac{1}{0.109}
      \approx 9.17.
\]

Несмотря на ускорение основной части в $100$ раз, вся программа ускорилась только примерно в $9.17$ раза.

\section{Закон Густафсона}

Закон Амдала фиксирует размер задачи. На практике при наличии более мощного вычислителя часто увеличивают размер задачи.

\begin{theorem}[Закон Густафсона]
Пусть доля последовательной части программы при выполнении на параллельной системе равна $\alpha$, а число вычислительных элементов равно $N$. Тогда масштабированное ускорение оценивается формулой:
\[
    S_G = N - \alpha(N-1).
\]
\end{theorem}

\begin{proof}
Нормируем время выполнения программы на параллельной системе к единице:
\[
    T_N = 1.
\]
Пусть последовательная часть занимает долю $\alpha$, а параллельная часть --- долю $1-\alpha$:
\[
    T_N = \alpha + (1-\alpha) = 1.
\]

Если ту же масштабированную задачу выполнять на одном вычислительном элементе, то последовательная часть останется равной $\alpha$, а параллельная часть будет выполняться в $N$ раз дольше:
\[
    T_1 = \alpha + N(1-\alpha).
\]

Масштабированное ускорение:
\[
    S_G = \frac{T_1}{T_N}
        = \alpha + N(1-\alpha).
\]

Раскроем скобки:
\[
    S_G = \alpha + N - N\alpha
        = N - \alpha(N-1).
\]
Теорема доказана.
\end{proof}

\section{Модель Roofline}

Для анализа производительности GPU часто используется модель Roofline.

\begin{definition}
\textbf{Модель Roofline} --- простая модель производительности, связывающая достижимую производительность программы с ее арифметической интенсивностью, пиковой вычислительной производительностью и пропускной способностью памяти.
\end{definition}

Пусть:
\[
    P_{\max}
\]
--- пиковая производительность GPU в операциях с плавающей точкой в секунду, а
\[
    B
\]
--- пропускная способность памяти в байтах в секунду.

Если арифметическая интенсивность программы равна $I$, то достижимая производительность ограничена:
\[
    P \le \min(P_{\max}, BI).
\]

\begin{proposition}[Ограничение Roofline]
Для программы с арифметической интенсивностью $I$ достижимая производительность не превосходит
\[
    P^* = \min(P_{\max}, BI).
\]
\end{proposition}

\begin{proof}
С одной стороны, производительность не может превышать аппаратный пик:
\[
    P \le P_{\max}.
\]

С другой стороны, если программа передает из памяти $M$ байт, а выполняет $F$ арифметических операций, то
\[
    I = \frac{F}{M}.
\]
При пропускной способности памяти $B$ минимальное время передачи данных:
\[
    T_{\text{mem}} \ge \frac{M}{B}.
\]
Следовательно, производительность по операциям:
\[
    P = \frac{F}{T} \le \frac{F}{T_{\text{mem}}}
      \le \frac{F}{M/B}
      = B\frac{F}{M}
      = BI.
\]

Итак,
\[
    P \le P_{\max}
\]
и
\[
    P \le BI.
\]
Следовательно,
\[
    P \le \min(P_{\max}, BI).
\]
Утверждение доказано.
\end{proof}

\subsection{Интерпретация}

Если
\[
    BI < P_{\max},
\]
то программа ограничена памятью.

Если
\[
    BI \ge P_{\max},
\]
то программа ограничена вычислительными ресурсами.

Граница между режимами определяется условием:
\[
    I = \frac{P_{\max}}{B}.
\]

\section{Программные модели для GPU}

\subsection{CUDA}

\begin{definition}
\textbf{CUDA} --- платформа и модель программирования NVIDIA для параллельных вычислений на GPU.
\end{definition}

CUDA предоставляет:

\begin{itemize}
    \item расширения языка C/C++;
    \item компилятор \texttt{nvcc};
    \item API управления памятью GPU;
    \item библиотеки линейной алгебры, БПФ, случайных чисел и глубокого обучения;
    \item средства профилирования и отладки.
\end{itemize}

Пример CUDA-ядра для сложения векторов:

\begin{lstlisting}[style=cuda]
__global__ void vecAdd(const float* a, const float* b, float* c, int n) {
    int i = blockIdx.x * blockDim.x + threadIdx.x;
    if (i < n) {
        c[i] = a[i] + b[i];
    }
}
\end{lstlisting}

Запуск ядра:

\begin{lstlisting}[style=cuda]
int threads = 256;
int blocks = (n + threads - 1) / threads;
vecAdd<<<blocks, threads>>>(d_a, d_b, d_c, n);
\end{lstlisting}

Здесь:
\[
    \texttt{threads} = 256
\]
--- число потоков в блоке, а
\[
    \texttt{blocks} = \left\lceil \frac{n}{256} \right\rceil
\]
--- число блоков.

\subsection{OpenCL}

\begin{definition}
\textbf{OpenCL} --- открытый стандарт для параллельного программирования гетерогенных систем, включающих CPU, GPU, FPGA и другие ускорители.
\end{definition}

OpenCL менее привязан к конкретному производителю, чем CUDA. Он позволяет писать переносимый код, но часто требует более подробного управления платформами, устройствами, контекстами, очередями команд и буферами.

\subsection{SYCL}

\begin{definition}
\textbf{SYCL} --- высокоуровневая C++-модель программирования для гетерогенных систем, ориентированная на переносимость между различными ускорителями.
\end{definition}

SYCL позволяет писать однопроходный C++-код, где host- и device-части теснее интегрированы, чем в классическом OpenCL.

\subsection{Библиотеки}

На практике часто выгоднее использовать готовые GPU-библиотеки:

\begin{itemize}
    \item cuBLAS --- линейная алгебра на GPU;
    \item cuDNN --- примитивы для нейронных сетей;
    \item cuFFT --- быстрое преобразование Фурье;
    \item Thrust --- параллельные алгоритмы в стиле STL;
    \item NCCL --- коммуникации между GPU;
    \item rocBLAS, MIOpen --- аналоги для экосистемы AMD ROCm.
\end{itemize}

\section{Алгоритм сложения векторов на GPU}

\subsection{Постановка задачи}

Даны два массива:
\[
    a = (a_0,a_1,\dots,a_{n-1}),
\]
\[
    b = (b_0,b_1,\dots,b_{n-1}).
\]
Нужно вычислить:
\[
    c_i = a_i + b_i,\qquad i = 0,1,\dots,n-1.
\]

\subsection{Шаги алгоритма}

\begin{enumerate}
    \item Выделить память на CPU для массивов $a$, $b$, $c$.
    \item Инициализировать входные массивы $a$ и $b$.
    \item Выделить память на GPU для массивов $d_a$, $d_b$, $d_c$.
    \item Скопировать $a$ и $b$ из памяти CPU в память GPU.
    \item Запустить CUDA-ядро, где каждый поток вычисляет один элемент $c_i$.
    \item Скопировать результат $d_c$ обратно в память CPU.
    \item Освободить память.
\end{enumerate}

\subsection{Иллюстрация распределения работы}

\[
\begin{array}{c|cccccccc}
\text{Поток} & 0 & 1 & 2 & 3 & 4 & 5 & 6 & 7 \\
\hline
\text{Операция} &
c_0=a_0+b_0 &
c_1=a_1+b_1 &
c_2=a_2+b_2 &
c_3=a_3+b_3 &
c_4=a_4+b_4 &
c_5=a_5+b_5 &
c_6=a_6+b_6 &
c_7=a_7+b_7
\end{array}
\]

\subsection{Минимальный пример}

Пусть:
\[
    a = (1,2,3,4),
    \qquad
    b = (10,20,30,40).
\]
Тогда:
\[
    c = (11,22,33,44).
\]

Каждый поток выполняет одну операцию:
\[
\begin{array}{c|c}
\text{Поток} & \text{Результат} \\
\hline
0 & c_0 = 1 + 10 = 11 \\
1 & c_1 = 2 + 20 = 22 \\
2 & c_2 = 3 + 30 = 33 \\
3 & c_3 = 4 + 40 = 44
\end{array}
\]

\section{Алгоритм редукции суммы}

\subsection{Постановка задачи}

Дан массив:
\[
    x = (x_0,x_1,\dots,x_{n-1}).
\]
Нужно вычислить:
\[
    S = \sum_{i=0}^{n-1} x_i.
\]

Редукция является типичной параллельной операцией. Она используется при вычислении скалярных произведений, норм векторов, статистик и агрегатов.

\subsection{Идея параллельной редукции}

Суммирование можно организовать деревом.

Например, для $8$ элементов:

\[
\begin{array}{cccccccc}
x_0 & x_1 & x_2 & x_3 & x_4 & x_5 & x_6 & x_7 \\
\downarrow & \downarrow & \downarrow & \downarrow & \downarrow & \downarrow & \downarrow & \downarrow \\
x_0+x_1 & & x_2+x_3 & & x_4+x_5 & & x_6+x_7 & \\
\downarrow & & \downarrow & & \downarrow & & \downarrow & \\
x_0+x_1+x_2+x_3 & & & & x_4+x_5+x_6+x_7 & & & \\
\downarrow & & & & \downarrow & & & \\
x_0+x_1+\cdots+x_7 & & & & & & &
\end{array}
\]

\subsection{Шаги алгоритма редукции внутри блока}

\begin{enumerate}
    \item Каждый поток загружает один или несколько элементов из глобальной памяти.
    \item Значения помещаются в разделяемую память.
    \item Выполняется синхронизация потоков блока.
    \item На каждом шаге расстояние между суммируемыми элементами уменьшается вдвое.
    \item После каждого шага выполняется синхронизация.
    \item Поток с индексом $0$ записывает частичную сумму блока в глобальную память.
    \item Частичные суммы блоков редуцируются повторным запуском ядра или на CPU.
\end{enumerate}

\subsection{Псевдокод}

\begin{lstlisting}[style=cuda]
__global__ void reduceSum(float* input, float* partial, int n) {
    extern __shared__ float s[];

    int tid = threadIdx.x;
    int i = blockIdx.x * blockDim.x + threadIdx.x;

    if (i < n)
        s[tid] = input[i];
    else
        s[tid] = 0.0f;

    __syncthreads();

    for (int stride = blockDim.x / 2; stride > 0; stride /= 2) {
        if (tid < stride) {
            s[tid] += s[tid + stride];
        }
        __syncthreads();
    }

    if (tid == 0) {
        partial[blockIdx.x] = s[0];
    }
}
\end{lstlisting}

\subsection{Минимальный пример}

Пусть:
\[
    x = (1,2,3,4,5,6,7,8).
\]

Шаг 1:
\[
    (1+5,\;2+6,\;3+7,\;4+8) = (6,8,10,12).
\]

Шаг 2:
\[
    (6+10,\;8+12) = (16,20).
\]

Шаг 3:
\[
    16+20 = 36.
\]

Итог:
\[
    S = 36.
\]

\section{Алгоритм умножения матриц на GPU}

\subsection{Постановка задачи}

Даны матрицы:
\[
    A \in \mathbb{R}^{m \times k}, \qquad
    B \in \mathbb{R}^{k \times n}.
\]
Нужно вычислить:
\[
    C = AB,\qquad C \in \mathbb{R}^{m \times n}.
\]

Элемент результата:
\[
    C_{ij} = \sum_{\ell=0}^{k-1} A_{i\ell}B_{\ell j}.
\]

\subsection{Наивный параллельный алгоритм}

Каждый поток вычисляет один элемент $C_{ij}$.

\begin{enumerate}
    \item Назначить каждому потоку пару индексов $(i,j)$.
    \item Поток считывает строку $i$ матрицы $A$ и столбец $j$ матрицы $B$.
    \item Поток вычисляет скалярное произведение:
    \[
        C_{ij} = \sum_{\ell=0}^{k-1} A_{i\ell}B_{\ell j}.
    \]
    \item Поток записывает результат в глобальную память.
\end{enumerate}

\subsection{Минимальный пример}

Пусть:
\[
A =
\begin{pmatrix}
1 & 2 \\
3 & 4
\end{pmatrix},
\qquad
B =
\begin{pmatrix}
5 & 6 \\
7 & 8
\end{pmatrix}.
\]

Тогда:
\[
C_{00} = 1\cdot 5 + 2\cdot 7 = 19,
\]
\[
C_{01} = 1\cdot 6 + 2\cdot 8 = 22,
\]
\[
C_{10} = 3\cdot 5 + 4\cdot 7 = 43,
\]
\[
C_{11} = 3\cdot 6 + 4\cdot 8 = 50.
\]

Следовательно:
\[
C =
\begin{pmatrix}
19 & 22 \\
43 & 50
\end{pmatrix}.
\]

\subsection{Тайловый алгоритм умножения матриц}

Пусть размер тайла равен $T \times T$.

\begin{enumerate}
    \item Разбить матрицы $A$, $B$, $C$ на тайлы.
    \item Один блок потоков отвечает за вычисление одного тайла матрицы $C$.
    \item Для каждого значения $q$:
    \begin{enumerate}
        \item загрузить тайл $A_{iq}$ в разделяемую память;
        \item загрузить тайл $B_{qj}$ в разделяемую память;
        \item синхронизировать потоки блока;
        \item накопить вклад произведения этих тайлов в локальную переменную;
        \item синхронизировать потоки перед загрузкой следующей пары тайлов.
    \end{enumerate}
    \item Записать вычисленный элемент $C_{ij}$ в глобальную память.
\end{enumerate}

\subsection{Псевдокод CUDA}

\begin{lstlisting}[style=cuda]
#define T 16

__global__ void matMulTiled(const float* A, const float* B, float* C,
                            int M, int K, int N) {
    __shared__ float As[T][T];
    __shared__ float Bs[T][T];

    int row = blockIdx.y * T + threadIdx.y;
    int col = blockIdx.x * T + threadIdx.x;

    float sum = 0.0f;

    for (int q = 0; q < (K + T - 1) / T; ++q) {
        int aCol = q * T + threadIdx.x;
        int bRow = q * T + threadIdx.y;

        if (row < M && aCol < K)
            As[threadIdx.y][threadIdx.x] = A[row * K + aCol];
        else
            As[threadIdx.y][threadIdx.x] = 0.0f;

        if (bRow < K && col < N)
            Bs[threadIdx.y][threadIdx.x] = B[bRow * N + col];
        else
            Bs[threadIdx.y][threadIdx.x] = 0.0f;

        __syncthreads();

        for (int e = 0; e < T; ++e) {
            sum += As[threadIdx.y][e] * Bs[e][threadIdx.x];
        }

        __syncthreads();
    }

    if (row < M && col < N) {
        C[row * N + col] = sum;
    }
}
\end{lstlisting}

\section{Сканирование префиксных сумм}

\subsection{Постановка задачи}

Для массива:
\[
    x = (x_0,x_1,\dots,x_{n-1})
\]
нужно вычислить массив префиксных сумм:
\[
    y_i = \sum_{j=0}^{i} x_j.
\]

Например:
\[
    x = (3,1,7,0,4),
\]
тогда:
\[
    y = (3,4,11,11,15).
\]

\subsection{Параллельный алгоритм Blelloch scan}

Для простоты рассмотрим exclusive scan, где:
\[
    y_0 = 0,
\]
\[
    y_i = \sum_{j=0}^{i-1} x_j.
\]

Алгоритм состоит из двух фаз:

\begin{enumerate}
    \item \textbf{Фаза подъема} \emph{up-sweep}: строится дерево частичных сумм.
    \item \textbf{Фаза спуска} \emph{down-sweep}: частичные суммы распространяются вниз по дереву.
\end{enumerate}

\subsection{Шаги алгоритма}

Пусть $n$ --- степень двойки.

\begin{enumerate}
    \item Скопировать входной массив во временный массив $a$.
    \item Для $d = 0,1,\dots,\log_2 n - 1$ выполнить:
    \[
        a[(k+1)2^{d+1}-1] \gets
        a[(k+1)2^{d+1}-1] + a[(k+1)2^{d+1}-2^d-1].
    \]
    \item После фазы подъема последний элемент содержит сумму всех элементов.
    \item Установить последний элемент равным нулю.
    \item Для $d = \log_2 n - 1,\dots,0$ выполнить обмен и сложение:
    \[
        t \gets a[(k+1)2^{d+1}-2^d-1],
    \]
    \[
        a[(k+1)2^{d+1}-2^d-1] \gets a[(k+1)2^{d+1}-1],
    \]
    \[
        a[(k+1)2^{d+1}-1] \gets a[(k+1)2^{d+1}-1] + t.
    \]
\end{enumerate}

\subsection{Минимальный пример}

Пусть:
\[
    x = (3,1,7,0).
\]

Фаза подъема:
\[
    (3,1,7,0)
    \to
    (3,4,7,7)
    \to
    (3,4,7,11).
\]

Устанавливаем последний элемент в ноль:
\[
    (3,4,7,0).
\]

Фаза спуска:
\[
    (3,0,7,4)
    \to
    (0,3,4,11).
\]

Итоговый exclusive scan:
\[
    y = (0,3,4,11).
\]

Inclusive scan получается добавлением исходного массива:
\[
    (0,3,4,11) + (3,1,7,0) = (3,4,11,11).
\]

\section{Синхронизация и взаимодействие потоков}

Внутри блока CUDA доступна барьерная синхронизация:
\[
    \texttt{\_\_syncthreads()}.
\]

Она гарантирует, что все потоки блока достигли данной точки программы и что операции с разделяемой памятью, выполненные до барьера, видны потокам блока после барьера.

\begin{remark}
Потоки из разных блоков в рамках одного запуска ядра обычно не имеют простой глобальной синхронизации. Если нужна глобальная синхронизация, часто используют несколько запусков ядер.
\end{remark}

\section{Occupancy}

\begin{definition}
\textbf{Occupancy} --- отношение числа активных warp на мультипроцессоре к максимально возможному числу активных warp.
\end{definition}

На occupancy влияют:

\begin{itemize}
    \item число потоков в блоке;
    \item число регистров на поток;
    \item объем разделяемой памяти на блок;
    \item аппаратные ограничения конкретной архитектуры.
\end{itemize}

Высокий occupancy помогает скрывать задержки памяти: пока один warp ожидает данные, другой warp может исполняться. Однако максимальный occupancy не всегда означает максимальную производительность. Иногда использование большего числа регистров или разделяемой памяти снижает occupancy, но повышает эффективность вычислений.

\section{Передача данных между CPU и GPU}

Передача данных между CPU и GPU часто является узким местом. Поэтому важно:

\begin{enumerate}
    \item минимизировать число копирований;
    \item передавать данные крупными блоками;
    \item по возможности перекрывать передачу данных и вычисления;
    \item использовать pinned memory для ускорения копирования;
    \item оставлять данные на GPU между последовательными этапами вычислений.
\end{enumerate}

Общая схема:
\[
    T_{\text{total}} =
    T_{\text{copy to GPU}} +
    T_{\text{kernel}} +
    T_{\text{copy from GPU}}.
\]

GPU-ускорение выгодно, если
\[
    T_{\text{CPU}} >
    T_{\text{copy to GPU}} +
    T_{\text{kernel}} +
    T_{\text{copy from GPU}}.
\]

\section{Потоки CUDA и асинхронность}

\begin{definition}
\textbf{CUDA stream} --- последовательность операций, выполняемых на GPU в заданном порядке. Операции из разных потоков могут выполняться параллельно, если это допускается аппаратурой и зависимостями.
\end{definition}

Потоки позволяют:

\begin{itemize}
    \item перекрывать копирование данных и вычисления;
    \item запускать несколько независимых ядер;
    \item организовывать конвейерную обработку больших данных.
\end{itemize}

Типичная конвейерная схема:

\[
\begin{array}{c|ccc}
\text{Часть данных} & \text{Копирование на GPU} & \text{Вычисление} & \text{Копирование на CPU} \\
\hline
1 & t_1 & t_2 & t_3 \\
2 & t_2 & t_3 & t_4 \\
3 & t_3 & t_4 & t_5
\end{array}
\]

\section{Численная точность на GPU}

GPU поддерживают различные форматы чисел:

\begin{itemize}
    \item FP64 --- двойная точность;
    \item FP32 --- одинарная точность;
    \item FP16 --- половинная точность;
    \item BF16 --- формат, широко используемый в машинном обучении;
    \item INT8, INT4 --- целочисленные форматы для инференса нейросетей.
\end{itemize}

Важно учитывать:

\begin{enumerate}
    \item операции с плавающей точкой не являются ассоциативными:
    \[
        (a+b)+c \ne a+(b+c)
    \]
    в общем случае из-за округления;
    \item параллельные редукции могут давать немного другой результат, чем последовательное суммирование;
    \item использование FP16 и BF16 ускоряет вычисления, но может снижать точность;
    \item для научных расчетов может требоваться FP64.
\end{enumerate}

\section{Типичные оптимизации GPU-программ}

\subsection{Оптимизация доступа к памяти}

\begin{itemize}
    \item использовать коалесцированные чтения и записи;
    \item избегать случайного доступа к глобальной памяти;
    \item использовать разделяемую память для повторного использования данных;
    \item уменьшать число обращений к глобальной памяти;
    \item выравнивать данные.
\end{itemize}

\subsection{Оптимизация вычислений}

\begin{itemize}
    \item избегать дивергенции внутри warp;
    \item выбирать разумный размер блока;
    \item использовать быстрые математические функции, если допустима меньшая точность;
    \item по возможности увеличивать арифметическую интенсивность;
    \item использовать специализированные блоки, например tensor cores, если задача подходит.
\end{itemize}

\subsection{Использование библиотек}

Если задача является стандартной, например умножение матриц или БПФ, обычно лучше использовать оптимизированную библиотеку. Такие библиотеки учитывают особенности конкретных архитектур и часто значительно превосходят самописные реализации.

\section{Примеры прикладного использования GPU}

\subsection{Линейная алгебра}

GPU особенно эффективны в операциях линейной алгебры:

\begin{itemize}
    \item умножение матриц;
    \item LU-, QR-, SVD-разложения;
    \item решение систем линейных уравнений;
    \item вычисление собственных значений;
    \item итерационные методы.
\end{itemize}

Основная причина эффективности --- высокая регулярность вычислений и большая арифметическая интенсивность.

\subsection{Машинное обучение}

Современное глубокое обучение активно использует GPU, потому что обучение нейронных сетей сводится к большому числу матричных операций и сверток. Основные операции:

\[
    Y = WX + b,
\]
\[
    Y_{i,j,k} = \sum_{u,v,c} X_{i+u,j+v,c}K_{u,v,c,k}.
\]

Эти операции хорошо распараллеливаются.

\subsection{Обработка изображений}

Изображение можно рассматривать как двумерный массив пикселей. Многие фильтры применяются независимо или почти независимо к каждому пикселю:

\[
    y_{ij} = f(x_{ij}).
\]

Например, инверсия яркости:
\[
    y_{ij} = 255 - x_{ij}.
\]

Сверточный фильтр:
\[
    y_{ij} = \sum_{u=-r}^{r}\sum_{v=-r}^{r} k_{uv}x_{i+u,j+v}.
\]

\subsection{Метод Монте-Карло}

Методы Монте-Карло основаны на большом числе независимых случайных экспериментов. Это хорошо подходит для GPU.

Например, оценка числа $\pi$:

\begin{enumerate}
    \item Сгенерировать много случайных точек $(x,y)$ в квадрате $[0,1]\times[0,1]$.
    \item Проверить условие:
    \[
        x^2+y^2 \le 1.
    \]
    \item Подсчитать долю точек внутри четверти круга.
    \item Оценить:
    \[
        \pi \approx 4\frac{N_{\text{inside}}}{N}.
    \]
\end{enumerate}

Каждый поток GPU может обрабатывать одну или несколько случайных точек.

\subsection{Численное решение дифференциальных уравнений}

При решении уравнений в частных производных область разбивается на сетку. Значение в каждой точке обновляется по локальному шаблону.

Например, для одномерного уравнения теплопроводности:
\[
    u_i^{t+1} =
    u_i^t + \alpha
    \left(
        u_{i-1}^t - 2u_i^t + u_{i+1}^t
    \right).
\]

Каждый поток может обновлять одну точку сетки. Так как используются соседние значения, полезна разделяемая память.

\section{Когда GPU может быть неэффективен}

GPU не всегда дает ускорение. Неблагоприятные случаи:

\begin{itemize}
    \item малый размер задачи;
    \item сложная ветвящаяся логика;
    \item сильные зависимости между шагами;
    \item нерегулярные обращения к памяти;
    \item частые обмены CPU--GPU;
    \item низкая арифметическая интенсивность;
    \item необходимость большого количества синхронизаций;
    \item преобладание последовательной части алгоритма.
\end{itemize}

Пример плохой задачи: обработка небольшого массива из нескольких сотен элементов со сложными условными переходами. Накладные расходы запуска ядра и копирования данных могут превысить выигрыш от параллелизма.

\section{Общая методика переноса задачи на GPU}

\begin{enumerate}
    \item Выделить наиболее трудоемкие участки программы с помощью профилирования.
    \item Определить, есть ли в них массовый параллелизм.
    \item Оценить объем передаваемых данных и арифметическую интенсивность.
    \item Выбрать модель программирования: CUDA, OpenCL, SYCL или библиотеку.
    \item Спроектировать разбиение данных на потоки, блоки и сетку.
    \item Обеспечить коалесцированный доступ к памяти.
    \item Использовать разделяемую память, если есть повторное использование данных.
    \item Минимизировать дивергенцию.
    \item Проверить корректность результатов.
    \item Провести профилирование и оптимизацию.
\end{enumerate}

\section{Краткий итог}

Графические процессоры являются мощными ускорителями для задач с массовым параллелизмом. Их архитектура ориентирована на высокую пропускную способность, а не на минимальную задержку одного потока. Наиболее эффективно GPU применяются в задачах линейной алгебры, машинного обучения, обработки изображений, численного моделирования и статистических вычислений.

Ключевые факторы эффективности GPU-программ:

\begin{itemize}
    \item большое число независимых операций;
    \item регулярный и коалесцированный доступ к памяти;
    \item высокая арифметическая интенсивность;
    \item низкая дивергенция потоков;
    \item эффективное использование разделяемой памяти;
    \item минимизация обменов между CPU и GPU;
    \item использование оптимизированных библиотек.
\end{itemize}

Главная практическая идея состоит в том, что GPU следует использовать не как быстрый универсальный процессор, а как специализированный массово-параллельный ускоритель для подходящих фрагментов вычислений.

\end{document}
